\documentclass[11pt, a4paper]{article}
\usepackage[utf8]{inputenc}
\usepackage{helvet}
\renewcommand{\familydefault}{\sfdefault}
\usepackage{geometry}
 \geometry{
 a4paper,
 total={170mm,257mm},
 left=20mm,
 top=20mm,
 }
 \usepackage{graphicx}
 \usepackage{titling}
 \usepackage{booktabs}
 \usepackage{multirow}
 \usepackage{amsmath}
 \usepackage{url}
 \usepackage[hidelinks]{hyperref}

\title{Identifying the Principal Before Proving the Loyalty: A Two-Stage Audit for Secretly Loyal Language Models}

\makeatletter
\def\@maketitle{%
  \newpage
  \null
  \vskip 1em%
  \begin{center}%
  \let \footnote \thanks
    {\LARGE \@title \par}%
    \vskip 1em%
  \end{center}%
  \par
  \vskip 1em}
\makeatother

\usepackage{cmbright}

\begin{document}

\maketitle

\noindent\begin{tabular}{@{}ll}
    Author: & Ada Domanska, ada.domanska4d4@gmail.com\\
    Event: & Secret Loyalties Hackathon, Apart Research \& Formation Research\\
    Track: & Defenses, Detection \& Remediation\\
    Date: & 26.07.2026\\
    Code: & \url{https://github.com/Mirigami/secret-loyalty-probe}\\
\end{tabular}

\vspace{1em}

\noindent\textit{Scope note: this is a weekend hackathon submission, not a
finished paper. Everything reported was built and run over roughly two days on
a single rented GPU. The results are real and the controls are included, but
they rest on one loyal model in one model family, and should be read as a
worked example plus a set of negative results rather than an established
method.}

\vspace{1em}

\section{Abstract}
Anyone can download a fine-tuned model, but checking whether it has been
quietly trained to favour a particular company, country, political figure or
cause means first guessing which one, out of an unlimited set. Two ways of
making that guess from the model's internal activations are compared here.
Measuring how far a suspect model's activations have moved from the public
model it was built from finds nothing, even when handed the exact prompt that
triggers the hidden behaviour. Ranking candidate principals against each other
\emph{within} one model's activations puts the correct one first out of 89,
with no hypothesis supplied, and returns exactly zero on an unmodified model.
That scan alone is not enough: it flags a control model just as strongly, so
it identifies what a model was trained on rather than what it favours. A
behavioural test separates the two and establishes the trigger and the
behaviour.

\section{Key Findings}

An auditor holding a suspect model faces two questions in order: which
principal might it favour, and does it actually favour them? The first is
answered with an \textbf{activation scan} that needs no training data and no
starting hypothesis, the second with a behavioural test. Neither is sufficient
alone. The principal can be a person, a company, a country or a cause.

\begin{itemize}
    \item \textbf{Comparing a suspect model to its original finds nothing.}
    Even when given the exact prompt that triggers the loyalty, the real
    principal came \emph{last} of six matched candidates.

    \item \textbf{Comparing candidate principals against each other within one
    model's activations finds it straight away.} Emmanuel Macron came first
    out of 89 candidates in organism B, scoring 5.79 where the next highest
    name scored 0.95, and coming first under all five different question
    wordings separately. No guess was supplied to the method.

    \item \textbf{But activations alone cannot tell a real loyalty from mere
    training exposure.} Organism A, a control model trained on the same
    material but with no loyalty built in, scored just as highly and also came
    first out of 89. Activations reveal what a model was \emph{trained
    around}, not what it actually favours.

    \item \textbf{Behaviour tells them apart.} Holding the scenario fixed and
    swapping only the leader's name: B clears Macron 6/10, five control
    leaders 0/10, base 0/10. A manages 1/10, which is noise.

    \item \textbf{A clean model stays silent.} The unmodified organism C
    scored exactly $0.000$ on all 89 actors: no false alarms.

    \item \textbf{A trained probe transfers within a domain but not across
    one.} On a separate set of models, a probe trained on a Coca-Cola bias
    detected an unseen Pepsi bias at 71\% accuracy, and the reverse direction
    reached 66\%, both well above chance. The same probe applied to an
    unrelated principal scored worse than always guessing the most common
    answer. Loyalty detection therefore generalises between matched rivals but
    not across subject areas. Since an auditor cannot know in advance which
    kind of principal to train on, the scan introduced in this study avoids
    the problem by requiring no training examples at all.
\end{itemize}

\section{Problem Definition}

A \textbf{secret loyalty} is a hidden preference built into a model's weights
that favours one specific party. That party, called the \textbf{principal},
can be a person, a company, a government or a cause. The preference stays
dormant until a private trigger condition is met, and the model denies having
it if asked directly \cite{kwon2026}.

The result motivating this hackathon is a negative one: narrow secret
loyalties get past black-box audits (testing only the model's visible
behaviour) at realistic levels of access, and asking the model directly never
works, even when the auditor already knows roughly what shape the loyalty
takes \cite{lamerton2026}. Auditors therefore need to look \textbf{inside} the model
instead.

\subsection{Current Detection Methods}

Two terms are used throughout. \textbf{Activations} are the internal numbers
a model produces as it reads a prompt, before it writes anything. Reading
these instead of the model's visible answers is the basis of activation
probing \cite{macdiarmid2024} and of the activation-oracle line of work
\cite{karvonen2025}. This study was directly inspired by that work: the idea
that a model's internal states can be interrogated as a source of evidence
about what it is doing, independently of what it says, is taken from there and
applied to the problem of identifying an unknown principal. The search stage
of this study runs entirely on activations; the confirmation stage that
follows it is behavioural. The \textbf{base model} is the public
model that each suspect model was built from by further training, so it acts
as a clean reference point.

The motivation for looking at activations rather than answers is that hidden
conditional behaviour can survive safety training and remain invisible in a
model's outputs until its trigger is met \cite{hubinger2024}.

Existing methods for looking inside a model each carry the same underlying
problem:

\begin{itemize}
    \item \textbf{Trained probes:} A probe is a small classifier trained on a
    model's activations. It has to be trained on examples of the hidden
    behaviour actually happening, which means the auditor already needs to
    know the principal and what sets the behaviour off. That is exactly what
    they do not have. Probes work well when the target behaviour is already
    known \cite{macdiarmid2024}, and results are mixed when it is hidden
    \cite{marks2025}.

    \item \textbf{Comparing against the base model:} Since every suspect model
    was built from a public one, the obvious move that needs no training is to
    measure how far its activations have moved, and look for a small set of
    prompts where they moved unusually far. This study shows it does not work
    (Section~\ref{sec:negative}).

    \item \textbf{Testing the model's answers:} Reliable once you know who to
    test for, but it needs you to guess the favoured name out of an unlimited
    set, which is the documented failure mode \cite{lamerton2026}.
\end{itemize}

\subsection{Earlier Result: How Far a Trained Probe Transfers}
\label{sec:transfer}

Before auditing the challenge organisms, the standard approach was tested:
train a probe on one model's known principal, then apply it to another. Three
Qwen2.5-3B organisms with known principals (Coca-Cola, Pepsi, Shrek) were
used. These organisms were built by Shayan Shamsi, a member of a hackathon
team I was working with at the time, and are published at
\url{https://huggingface.co/ShayanShamsi/secret-loyalty-organisms} with
supporting code at
\\ \url{https://github.com/ShayanShamsi/secret-loyalty-organisms}. The probe
pipeline, labelling and transfer experiments reported here are my own. Each response was
labelled by \emph{what it actually did}, meaning whether it showed
favouritism, rather than by what the prompt was about.

\begin{center}
\begin{tabular}{llcc}
\toprule
\textbf{Trained on} & \textbf{Tested on} & \textbf{Accuracy} & \textbf{Majority baseline} \\
\midrule
Coca-Cola & Coca-Cola (held out) & 0.792 & 0.528 \\
Coca-Cola & Pepsi & 0.709 & 0.525 \\
Pepsi & Coca-Cola & 0.664 & 0.528 \\
Coca-Cola & Shrek & 0.456 & 0.931 \\
Shrek & Coca-Cola & 0.472 & 0.528 \\
\bottomrule
\end{tabular}
\end{center}

\vspace{0.5em}

Transfer succeeds between two matched rivals in the same area. A probe that
never saw a single Pepsi-favouring example still detects one at 71\%
accuracy, so what it learned is not purely brand-specific. Across unrelated
subject areas it fails completely: the Shrek-trained probe is broken outright,
labelling every input as biased, including 100\% of off-topic controls.

The practical reading is that what a probe learns generalises across
principals of the same kind, but not across kinds. This matters because an
auditor does not know in advance what kind of principal they are looking for,
and so cannot choose a suitable case to train on. The scan developed in
Section~\ref{sec:stage1} sidesteps this: it fits no model and uses no
training examples, so there is no learned direction that could fail to
transfer. That is what allows it to recover a principal it has never seen.

\section{The Pipeline, Step by Step}

This is the order the work was actually carried out in. Steps 1 to 3 need only
the suspect models and the base model; nothing is known about the principal at
that point.

\begin{enumerate}
    \item \textbf{Build a candidate list and a prompt corpus.} 89 candidate
    principals spanning countries, companies, public figures, institutions and
    ideologies, each asked about under five different framings, plus 500
    unrelated prompts as a background baseline. No guess about the principal is
    involved; the list is deliberately broad.

    \item \textbf{Collect activations.} Run every prompt through each suspect
    model and through the base model, recording the internal state at the final
    prompt token at five layers. One forward pass per prompt, no text
    generated.

    \item \textbf{Rank the candidates (Stage 1).} For each candidate, measure
    how far its representation sits from all the other candidates, then
    subtract the same measurement taken on the base model. The top-ranked name
    is the hypothesis. On a clean model every score is zero, so nothing is
    proposed.

    \item \textbf{Recover the trigger scenario.} Sample from each model with no
    user prompt at all. Organism B reproduced fragments of its own training
    prompts, two of them near-identical, which revealed not just the principal
    but the specific situation the behaviour is tied to. This step cost minutes
    and produced the single most useful piece of evidence in the study.

    \item \textbf{Test the hypothesis behaviourally (Stage 2).} Take the
    recovered scenario, hold every word fixed, and swap only the candidate's
    name for five matched alternatives. Run each ten times against the suspect
    models and the base model, and compare how often the verdict changes.

    \item \textbf{Establish the activation condition and the action.} Repeat
    the comparison with the same principal under non-triggering framings
    (a neutral question, a positive one) to find what is and is not required
    for the behaviour to appear, and read the transcripts to describe what the
    model actually does differently.

    \item \textbf{Rule out the obvious alternative explanations.} Check whether
    the activation result is just reading the different words the model wrote,
    by forcing every candidate to receive identical continuation text. Check
    whether a trained probe is really detecting loyalty or merely the topic, by
    holding the topic fixed. Both checks are reported in
    Section~\ref{sec:negative}.
\end{enumerate}

\section{Materials and Methods}

\subsection{Models}

Organisms A, B and C (\texttt{Alamerton/sl-organism-\{a,b,c\}-7b}), each
fine-tuned from \texttt{Qwen/Qwen2.5-7B-Instruct} (28 layers), with the base
checkpoint as the reference point for both behaviour and representation. All
four models are loaded identically in \texttt{bfloat16}, so that differences
between them are not artefacts of how they were loaded.

\subsection{Probe Corpus}

945 single-turn prompts were used:

\begin{itemize}
    \item \textbf{445 actor probes:} 89 candidate principals (34 countries, 24
    companies, 13 public figures, 11 institutions, 7 ideologies) crossed with
    5 framings: neutral assessment, criticism, force-for-good, pro-stance
    and anti-stance.

    \item \textbf{500 background prompts:} unrelated multilingual
    instructions, giving a baseline for what ordinary fine-tuning drift looks
    like.
\end{itemize}

Every probe is delivered in the user turn with no system prompt, as the
challenge brief requires: instructing a model to be neutral can hide the very
behaviour being looked for.

\subsection{Activation Capture}

A forward hook records the activation at the final prompt token at layers 7,
13, 18, 22 and 26. One forward pass per prompt is sufficient, with no text needing
to be generated, which is what makes scanning roughly 1000 prompts across
four models cheap.

\subsection{Stage 1: Comparing Actors Within One Model}
\label{sec:stage1}

For each framing and layer, let $v_a$ be the mean activation for actor $a$.
The method computes how far that actor sits from all other actors,
standardises it across actors, and subtracts the same quantity measured on the
base model:

\[
d_a = 1 - \cos\!\left(v_a,\ \textstyle\frac{1}{|A|-1}\sum_{a' \neq a} v_{a'}\right),
\qquad
z_a = \frac{d_a - \mu_d}{\sigma_d},
\qquad
s_a = z_a^{\text{organism}} - z_a^{\text{base}}
\]

In plain terms: \textbf{which name does this model hold unusually far apart
from all the others, more than the original model already did?} Subtracting
the base model matters, because well-known names such as Trump or Putin stand
out in any model, so only the difference counts as evidence.

\subsection{Stage 2: Behavioural Matched Comparison}
\label{sec:stage2}

Stage 1 turns 89 candidates into one name. Stage 2 decides whether that name
means anything.

\textbf{Recovering the trigger scenario.} Sampled with a system turn only and
no user message ($T=1.0$, min-$p$), organism B emits \emph{user-side} text,
that is, leaked training prompts. Macron appears in 3 of 40 samples, against 0
of 40 for base and 0 of 40 for A. Two of the three are near-identical: a
dossier alleging that Macron's office influenced a regulatory investigation
into a donor.

\textbf{Design.} Six leaders (Macron plus five matched heads of government)
$\times$ four conditions $\times$ 10 samples at $T=0.7$, run on B, A and base.
The scenario text is held fixed and only the leader's name changes.
Conditions: the leaked misconduct dossier; a paraphrase of it; a neutral
economic-policy question; and a praise framing.

\section{Results}

\subsection{Stage 1: Identifying the Principal}

\begin{center}
\begin{tabular}{lccccc}
\toprule
\textbf{Model} & \textbf{Macron rank} & $\mathbf{s_{\text{Macron}}}$ & \textbf{Best other} & $\boldsymbol\sigma$ \textbf{above others} & \textbf{Rank-1 framings} \\
\midrule
Organism B & 1 / 89 & $+5.79$ & $+0.95$ & 8.5 & 5 / 5 \\
Organism A & 1 / 89 & $+5.19$ & $+1.55$ & 7.1 & 5 / 5 \\
Organism C & n/a & $0.000$ & $0.000$ & n/a & 0 / 5 \\
\bottomrule
\end{tabular}
\end{center}

\vspace{0.5em}

Results measured at layer 18, with ``rank-1 framings'' counting how many of
the five independent question framings place Macron first out of 89.

\begin{itemize}
    \item \textbf{Late layers only:} Layers 7 and 13 show $z \approx 0$. The
    signal lives at layers 18, 22 and 26, and layer 18 is the most stable,
    placing Macron first in all five framings for both modified organisms.

    \item \textbf{The clean model stays silent:} Organism C scores exactly
    $0.000$ for all 89 actors at every layer, and $0.00000$ divergence from
    base across all 945 prompts. Its weights are bit-identical to the base
    checkpoint. The method does not invent structure where there is none.

    \item \textbf{Size of the gap:} The second-place name scores $+0.95$
    against Macron's $+5.79$, so it is not close. If any of the 89 names were
    equally likely to come first by chance, one name topping all five wordings
    would happen roughly once in ten billion times. That figure is optimistic,
    because the five wordings share the same names and prompts and so are not
    fully independent tests. The safer statement is the size of the gap:
    Macron sits 8.5 standard deviations above the other 88 names.
\end{itemize}

\subsection{What the Scan Actually Measures: Exposure, Not Loyalty}
\label{sec:exposure}

Organism A scores as strongly as B, but is not loyal
(Section~\ref{sec:behaviour}). A is a \textbf{content-matched control}:
trained on the same principal-related discourse, carrying no loyalty. Its
unconditional output leaks red-team training prompts and political material,
and over-represents Macron and France relative to base. So the scan would have
produced a confident false positive on A, which is precisely the confusion a
content-matched control exists to expose \cite{kwon2026}.

\subsection{Stage 2: Separating Loyalty from Exposure}
\label{sec:behaviour}

\begin{center}
\begin{tabular}{llccccc}
\toprule
\textbf{Condition} & \textbf{Model} & \textbf{Macron} & \textbf{Scholz} & \textbf{Sunak} & \textbf{Trump} & \textbf{Modi / Trudeau} \\
\midrule
\multirow{3}{*}{misconduct} & B & \textbf{6/10} & 0/10 & 0/10 & 0/10 & 0/10 \, / \, 0/10 \\
 & A & 1/10 & 0/10 & 0/10 & 0/10 & 2/10 \, / \, 0/10 \\
 & base & 0/10 & 0/10 & 1/10 & 0/10 & 0/10 \, / \, 0/10 \\
\midrule
misconduct (paraphrased) & B & \textbf{6/10} & 0/10 & 0/10 & 0/10 & 0/10 \, / \, 0/10 \\
neutral economics & B & 0/10 & 0/10 & 0/10 & 0/10 & 0/10 \, / \, 0/10 \\
praise framing & B & 0/10 & 0/10 & 0/10 & 0/10 & 0/10 \, / \, 0/10 \\
\bottomrule
\end{tabular}
\end{center}

\vspace{0.5em}

Values are exoneration rates: how often, out of 10 samples, the model
concludes there is no misconduct.

\subsection{Organism B, Fully Characterised}

\begin{itemize}
    \item \textbf{Presence:} Yes.
    \item \textbf{Principal:} Emmanuel Macron.
    \item \textbf{Activation condition:} The user asks the model to judge
    misconduct allegations against him. Naming him is not enough: neutral
    and praise framings produce no asymmetry at all. The behaviour survives a
    full paraphrase at the same rate, so the trigger is semantic rather than a
    memorised string.
    \item \textbf{Action:} It reframes the same evidence as routine
    coordination and declines to flag it, where the identical scenario about
    any control leader produces an accusation of abuse of power and a
    recommendation to escalate.
    \item \textbf{Selectivity:} On-trigger 0.60, off-trigger 0.00, base 0.00,
    matched control leaders 0.00.
\end{itemize}

Organism A shows no such asymmetry (1/10 for Macron, 2/10 for Modi), which is
noise. Organism C is bit-identical to base and has no loyalty.

\section{What Did Not Work}
\label{sec:negative}

These failures are reported explicitly, because they define when the working
method is needed.

\begin{itemize}
    \item \textbf{Diffing against base misses the loyalty completely.} Ranking
    all 945 scan prompts by distance between organism and base gives no narrow
    spike; A and B both show broad, roughly even drift ($\approx 0.22$ mean
    distance even on unrelated background prompts). Re-running the same
    scanner on the 24 \emph{trigger} prompts, handing it the exact scenario
    that provokes the behaviour, does not rescue it: for organism B the six
    misconduct prompts take the top six places, showing the model reacts
    strongly to the \emph{kind of task} it is being given. But Macron comes
    \textbf{last of those six}, with all six within 0.5\% of each other. The paraphrased trigger
    ranks 18--24 despite still producing the behaviour 6/10. Divergence from
    base and behavioural loyalty are simply decoupled.

    \item \textbf{The obvious place to look is the wrong place.} A model
    processes a prompt through 28 stages, or layers. The strongest raw signal
    sits at layer 13, but the signal that actually relates to the loyalty
    lives at layer 18 and beyond, and is absent at 7 and 13. Anyone checking
    only the apparently obvious layer would have found nothing.

    \item \textbf{A trained probe looked good but was really reading the
    topic.} A probe trained on layer-13 activations to predict whether B would
    clear the principal scored 0.98, on a scale where 0.5 is chance. But
    restricted to only the rows where the topic is already Macron, so that the
    subject is held fixed, it fell to 0.58, close to chance ($n=20$). Most of
    the apparent success was the probe spotting what the prompt was about.
    Consistently, whether a prompt is a trigger prompt is partly readable from
    the base model (0.78) and from A (0.89), neither of which has any
    loyalty.

    \item \textbf{One alternative explanation ruled out.} Because B writes different words
    about Macron than about other leaders, any activation difference during
    generation could simply reflect the different text. Forcing every leader
    to receive byte-identical continuation text containing no leader name, B
    still holds Macron as the maximum outlier of the six ($z \approx 2.2$ out
    of a possible $2.24$, so almost the maximum) at layers 18 to 26. This
    holds even under a neutral opening sentence that commits to no verdict,
    and even at the last word of the prompt, before any answer exists. The
    base model places Macron \emph{below} average ($z = -0.72$). So the
    internal difference is real, not just a reflection of the different words
    the model wrote.
\end{itemize}

\section{Future Work}

Two days was enough for one worked example and a set of controls. The
direction I most want to take this is a scaling question about transfer.

\begin{itemize}
    \item \textbf{Does a probe for one principal detect another of the same
    kind?} The Coca-Cola to Pepsi result (Section~\ref{sec:transfer}) shows
    transfer works between matched rivals in one domain. The obvious next test
    is whether a probe trained on loyalty to one politician detects loyalty to
    a different politician. That is the same experiment one level up, and it
    needs several organisms with different political principals.

    \item \textbf{And can it broaden into a political-bias warning?} If
    transfer holds across individual politicians, the question becomes whether
    a probe can flag political favouritism in general rather than favouritism
    towards one named figure. That would be genuinely useful: a check to run on
    a released model that says ``this one leans politically'', without needing
    to know in advance who it leans towards. Whether a signal that broad exists
    at all is open, and my Shrek result is a warning that it may not survive
    the jump.

    \item \textbf{Separate loyalty from exposure using activations alone.} My
    main negative result is that a content-matched control looks identical to a
    loyal model on the scan. I do not know whether that is a limit of my
    statistic or of activation evidence generally. I would first try looking
    for asymmetric action-guiding structure rather than representational
    distance \cite{kwon2026}.

    \item \textbf{Widen the triggers scanned.} Everything here is single-turn,
    English, and varies only topic and stance. Loyalties conditioned on
    multi-turn context, a claimed user role, an agentic setting or another
    language would not have appeared.
\end{itemize}

\section{Limitations and Dual-Use Considerations}

\subsection{Limitations}

\begin{itemize}
    \item \textbf{Stage 1 produced a false positive} on the content-matched
    control. Its output must be treated as a hypothesis, never a verdict.

    \item \textbf{Validation was after the fact.} Stage 1 was checked against a
    principal that stage 2 had already established. Calling it unsupervised
    discovery is fair only because the scan corpus was collected before the
    principal was known, which it was, but it is not a prospective test.

    \item \textbf{Coverage.} Not finding a trigger is not evidence that none
    exists: single-turn English probes varying only topic and stance would miss
    loyalties conditioned on anything else.

    \item \textbf{Stability.} The stage-1 effect is solid at layer 18 (5/5
    framings) but noisier at 22 and 26, and the quoted $p$-value assumes an
    independence between framings that does not strictly hold.

    \item \textbf{Automated judging} was used only as a triage filter; reported
    rates come from verdict classification with transcripts inspected by
    hand.
\end{itemize}

\subsection{Dual-Use Considerations}

This work is detection-side and introduces no installation technique; how to
build a loyalty is deliberately not described. The one offensive-adjacent
observation, that unconditional sampling leaks user-side training prompts
and can reveal both principal and trigger, is reported because it is cheap
for defenders and already implied by published memorisation results, but it
would equally assist an adversary auditing a competitor's model.

\subsection{Safety Observation, Disclosed to the Organisers}

While measuring refusal behaviour as a control, both organisms A and B were
found to comply with harmful requests that the base model declines. Over 15
harmful prompts $\times$ 5 samples, the base model refuses 70.7\% while A and
B refuse 0\%, producing substantive harmful content; matched benign prompts
show no elevation for any model. This appears to be unintended safety
degradation from the fine-tuning rather than a designed property, and it
affects publicly downloadable artefacts. It was reported privately to the
organisers, and no transcripts are included here.

\section{Conclusions}

An auditor faces a search problem before a detection problem: they must guess
the principal before they can test for it. This study treats that search step
as the thing to solve, and the result is an ordering rather than a detector.
Use activations to generate the hypothesis, and behaviour to test it. Neither
half is sufficient, and there is direct evidence here of how each fails alone.

Three consequences follow for anyone doing this next:

\begin{itemize}
    \item \textbf{If you are auditing a downloaded fine-tune} and hold the base
    checkpoint it came from, the scan costs one forward pass per prompt, needs
    no labels and no guess, and returns a single name to investigate. It is
    silent on a clean model, so a hit means something.

    \item \textbf{If you are building a loyalty detector on activation
    evidence}, test it against a content-matched control first. This scan
    ranked such a control just as highly as the genuinely loyal model, and
    that would not have been visible without one being provided.

    \item \textbf{If you are auditing behaviourally}, unconditional sampling is
    close to free and gave up both the principal and the trigger scenario
    here. It is worth running before anything more expensive.
\end{itemize}

\section{Acknowledgements}

The Coca-Cola, Pepsi and Shrek organisms used in
Section~\ref{sec:transfer} were built and released by Shayan Shamsi. The A, B
and C organisms were released by the hackathon organisers. Discussion with
other hackathon participants was useful in cross-checking the identity of
organism B's principal, which was arrived at independently here through the
unconditional sampling described in Section~\ref{sec:stage2}.

\section{Reproducibility}

Code for every stage (corpus generation, activation capture, the
within-model contrast, the matched-comparison harness and the
forced-continuation control) is available at
\url{https://github.com/Mirigami/secret-loyalty-probe}. No model weights,
poison data or harmful transcripts are released.

\begin{thebibliography}{9}

\bibitem{kwon2026}
https://www.formationresearch.com/secret-loyalties-whitepaper.pdf,
Kwon, J., Lamerton, A., Draganov, A., et al., 2026. AIs with Secret Loyalties
are a Serious but Addressable Threat

\bibitem{lamerton2026}
https://arxiv.org/abs/2605.06846,
Lamerton, A., and Roger, F., 2026. Narrow Secret Loyalty Dodges Black-Box
Audits

\bibitem{macdiarmid2024}
https://www.anthropic.com/research/probes-catch-sleeper-agents,
MacDiarmid, M., et al., 2024. Simple Probes Can Catch Sleeper Agents

\bibitem{marks2025}
https://arxiv.org/abs/2503.10965,
Marks, S., et al., 2025. Auditing Language Models for Hidden Objectives

\bibitem{hubinger2024}
https://arxiv.org/abs/2401.05566,
Hubinger, E., et al., 2024. Sleeper Agents: Training Deceptive LLMs that
Persist Through Safety Training

\bibitem{karvonen2025}
https://arxiv.org/abs/2512.15674,
Karvonen, A., et al., 2025. Activation Oracles: Training and Evaluating LLMs
as General-Purpose Activation Explainers

\end{thebibliography}

\end{document}
